\subsection{Différence d'héritage net}
\label{sec:diff_heritage}

L'héritage net est donc :
\begin{align}
    H_{CTO} &= C_{CTO} - D_{CTO}, \\
    H_{AV} &= X - D_{AV}.
\end{align}
La différence entre les deux supports s'écrit :
\begin{equation}
    \Delta H := H_{CTO} - H_{AV} = (C_{CTO} - C_{AV}) + PS - s f(B_{CTO}) + g(B_{AV}),
\end{equation}
\text{où } \(s = \frac{C_{CTO}}{C_{CTO} + B}\) \text{ représente la part imputable du CTO dans le patrimoine}.

Cette identité permet de distinguer plusieurs contributions :
\begin{itemize}
    \item le surcroît de capital du CTO lié à l'absence de frais (\(C_{CTO} - C_{AV}\)),
    \item les prélèvements sociaux propres à l'AV (\(PS\)),
    \item les droits successoraux imputés au CTO (\(s f(B_{CTO})\)),
    \item et les droits spécifiques de l'AV (\(g(B_{AV})\)).
\end{itemize}

Il convient de souligner que, trop souvent dans le débat public, la comparaison entre les deux enveloppes est réduite à un simple contraste entre les droits successoraux imputés au CTO et ceux applicables à l'assurance-vie. Une telle approche, bien que séduisante par sa simplicité, conduit à des conclusions partielles et parfois erronées. En effet, elle occulte des déterminants essentiels tels que l'impact cumulé des frais de gestion, le poids des prélèvements sociaux sur l'assurance-vie, ou encore la neutralisation des plus-values latentes au décès dans le CTO. Toute analyse rigoureuse doit donc considérer l'ensemble des composantes de l'identité précédente, faute de quoi le jugement porté sur l'avantage relatif de l'un ou l'autre support reste incomplet.
